% GENERATED FILE. Edit canonical lessons, then run npm run generate:books.
% canonical-source-hash: fnv1a64:9ad98378e0bc837a
% canonical-lessons: SA-C249-varta, SA-C249-natakam, SA-C249-kavyam, SA-C249-citrapatah, SA-C249-khela, SA-R249-first-pass-words-from-chapters-241-244, SA-R249-second-pass-words-from-chapters-245-249

\chapter{News, Drama, Poem, Film, Game}
\label{ch:sa-news-drama-poem-film-game}
\hlchaptermodality{249}

\begin{chapteropening}
\textbf{By the end of this chapter:} \emph{I can say news, a drama, a poem, a film and a game.}

Complete the last lesson of chapter 249: I can say news, a drama, a poem, a film and a game.
\end{chapteropening}

\section[\texorpdfstring{\sk{वार्ता} (vārtā)}{वार्ता (vārtā)}]{\sk{वार्ता} (vārtā) — news}

\label{lesson:SA-C249-varta}



\begin{quote}
Before the new one: say the Sanskrit for a stake, then the Sanskrit for a fee.
\end{quote}

\subsection*{You'll want to know: \sk{वार्ता}}
\textbf{\sk{वार्ता}} — \emph{vārtā} — \textquotedblleft{}news\textquotedblright{}.

\begin{grammarlens}[title={What you've built}]
One more word for this chapter.
\end{grammarlens}

\subsection*{Your turn}
\begin{itemize}
\raggedright
  \item \emph{Say it:} \emph{vārtā}
  \item \emph{Say it:} \emph{vārtā}, once more
  \item \emph{Say it:} say \emph{vārtā}
  \item \emph{Recall:} say the Sanskrit for a stake, then the Sanskrit for a fee, then say \emph{vārtā} again
\end{itemize}

\begin{tcolorbox}[breakable,colback=teal!4,colframe=teal!35!black,title={Before you move on}]
What does \sk{वार्ता} mean? (\textquotedblleft{}News\textquotedblright{}.) Say it once more.
\end{tcolorbox}
\section[\texorpdfstring{\sk{नाटकम्} (nāṭakam)}{नाटकम् (nāṭakam)}]{\sk{नाटकम्} (nāṭakam) — a drama, a play}

\label{lesson:SA-C249-natakam}



\begin{quote}
Before the new one: say the Sanskrit for a fee, then the Sanskrit for news.
\end{quote}

\subsection*{You'll want to know: \sk{नाटकम्}}
\textbf{\sk{नाटकम्}} — \emph{nāṭakam} — \textquotedblleft{}a drama, a play\textquotedblright{}.

\begin{grammarlens}[title={What you've built}]
One more word for this chapter.
\end{grammarlens}

\subsection*{Your turn}
\begin{itemize}
\raggedright
  \item \emph{Say it:} \emph{nāṭakam}
  \item \emph{Say it:} \emph{nāṭakam}, once more
  \item \emph{Say it:} say \emph{nāṭakam}
  \item \emph{Recall:} say the Sanskrit for a fee, then the Sanskrit for news, then say \emph{nāṭakam} again
\end{itemize}

\begin{tcolorbox}[breakable,colback=teal!4,colframe=teal!35!black,title={Before you move on}]
What does \sk{नाटकम्} mean? (\textquotedblleft{}A drama, a play\textquotedblright{}.) Say it once more.
\end{tcolorbox}
\section[\texorpdfstring{\sk{काव्यम्} (kāvyam)}{काव्यम् (kāvyam)}]{\sk{काव्यम्} (kāvyam) — a poem, poetry}

\label{lesson:SA-C249-kavyam}



\begin{quote}
Before the new one: say the Sanskrit for news, then the Sanskrit for a drama.
\end{quote}

\subsection*{You'll want to know: \sk{काव्यम्}}
\textbf{\sk{काव्यम्}} — \emph{kāvyam} — \textquotedblleft{}a poem, poetry\textquotedblright{}.

\begin{grammarlens}[title={What you've built}]
One more word for this chapter.
\end{grammarlens}

\subsection*{Your turn}
\begin{itemize}
\raggedright
  \item \emph{Say it:} \emph{kāvyam}
  \item \emph{Say it:} \emph{kāvyam}, once more
  \item \emph{Say it:} say \emph{kāvyam}
  \item \emph{Recall:} say the Sanskrit for news, then the Sanskrit for a drama, then say \emph{kāvyam} again
\end{itemize}

\begin{tcolorbox}[breakable,colback=teal!4,colframe=teal!35!black,title={Before you move on}]
What does \sk{काव्यम्} mean? (\textquotedblleft{}A poem, poetry\textquotedblright{}.) Say it once more.
\end{tcolorbox}
\section[\texorpdfstring{\sk{चित्रपटः} (citrapaṭaḥ)}{चित्रपटः (citrapaṭaḥ)}]{\sk{चित्रपटः} (citrapaṭaḥ) — a film, a painted cloth}

\label{lesson:SA-C249-citrapatah}



\begin{quote}
Before the new one: say the Sanskrit for a drama, then the Sanskrit for a poem.
\end{quote}

\subsection*{You'll want to know: \sk{चित्रपटः}}
\textbf{\sk{चित्रपटः}} — \emph{citrapaṭaḥ} — \textquotedblleft{}a film, a painted cloth\textquotedblright{}.

\begin{grammarlens}[title={What you've built}]
One more word for this chapter.
\end{grammarlens}

\subsection*{Your turn}
\begin{itemize}
\raggedright
  \item \emph{Say it:} \emph{citrapaṭaḥ}
  \item \emph{Say it:} \emph{citrapaṭaḥ}, once more
  \item \emph{Say it:} say \emph{citrapaṭaḥ}
  \item \emph{Recall:} say the Sanskrit for a drama, then the Sanskrit for a poem, then say \emph{citrapaṭaḥ} again
\end{itemize}

\begin{tcolorbox}[breakable,colback=teal!4,colframe=teal!35!black,title={Before you move on}]
What does \sk{चित्रपटः} mean? (\textquotedblleft{}A film, a painted cloth\textquotedblright{}.) Say it once more.
\end{tcolorbox}
\section[\texorpdfstring{\sk{खेला} (khelā)}{खेला (khelā)}]{\sk{खेला} (khelā) — a game}

\label{lesson:SA-C249-khela}



\begin{quote}
Before the new one: say the Sanskrit for a poem, then the Sanskrit for a film.
\end{quote}

\subsection*{You'll want to know: \sk{खेला}}
\textbf{\sk{खेला}} — \emph{khelā} — \textquotedblleft{}a game\textquotedblright{}.

\begin{grammarlens}[title={What you've built}]
One more word for this chapter.
\end{grammarlens}

\subsection*{Your turn}
\begin{itemize}
\raggedright
  \item \emph{Say it:} \emph{khelā}
  \item \emph{Say it:} \emph{khelā}, once more
  \item \emph{Say it:} say \emph{khelā}
  \item \emph{Recall:} say the Sanskrit for a poem, then the Sanskrit for a film, then say \emph{khelā} again
\end{itemize}

\begin{tcolorbox}[breakable,colback=teal!4,colframe=teal!35!black,title={Before you move on}]
What does \sk{खेला} mean? (\textquotedblleft{}A game\textquotedblright{}.) Say it once more.
\end{tcolorbox}
\section[(dialogue)]{Words from chapters 241-244 — a first pass}

\label{lesson:SA-R249-first-pass-words-from-chapters-241-244}



\begin{quote}
Say the words for a film and a game.
\end{quote}

\subsection*{Your turn}
Say these aloud:

\begin{itemize}
\raggedright
  \item for the sake of, by which, by that, that, however
  \item easy to get, hard to get, complicated, easy to do, hard to do
  \item mighty, sick, healthy, happy, unhappy
  \item proper, improper, necessary, other, alone
\end{itemize}

\begin{tcolorbox}[breakable,colback=teal!4,colframe=teal!35!black,title={Before you move on}]
For the sake of? (\textbf{\sk{कृते}}, \emph{kṛte}.) By which? (\textbf{\sk{येन}}, \emph{yena}.) By that? (\textbf{\sk{तेन}}, \emph{tena}.) That? (\textbf{\sk{यत्}}, \emph{yat}.) However? (\textbf{\sk{परन्तु}}, \emph{parantu}.) Easy to get? (\textbf{\sk{सुलभः}}, \emph{sulabhaḥ}.) Hard to get? (\textbf{\sk{दुर्लभः}}, \emph{durlabhaḥ}.) Complicated? (\textbf{\sk{जटिलः}}, \emph{jaṭilaḥ}.) Easy to do? (\textbf{\sk{सुकरः}}, \emph{sukaraḥ}.) Hard to do? (\textbf{\sk{दुष्करः}}, \emph{duṣkaraḥ}.) Mighty? (\textbf{\sk{बलवान्}}, \emph{balavān}.) Sick? (\textbf{\sk{रुग्णः}}, \emph{rugṇaḥ}.) Healthy? (\textbf{\sk{स्वस्थः}}, \emph{svasthaḥ}.) Happy? (\textbf{\sk{सुखी}}, \emph{sukhī}.) Unhappy? (\textbf{\sk{दुःखी}}, \emph{duḥkhī}.) Proper? (\textbf{\sk{उचितः}}, \emph{ucitaḥ}.) Improper? (\textbf{\sk{अनुचितः}}, \emph{anucitaḥ}.) Necessary? (\textbf{\sk{आवश्यकः}}, \emph{āvaśyakaḥ}.) Other? (\textbf{\sk{अन्यः}}, \emph{anyaḥ}.) Alone? (\textbf{\sk{एकाकी}}, \emph{ekākī}.)
\end{tcolorbox}
\section[(dialogue)]{Words from chapters 245-249 — a second pass}

\label{lesson:SA-R249-second-pass-words-from-chapters-245-249}



\begin{quote}
Say the words for a tube and buttermilk.
\end{quote}

\subsection*{Your turn}
Say these aloud:

\begin{itemize}
\raggedright
  \item a tube, buttermilk, fruit juice, a clove, cardamom
  \item puffed rice, grain, barley flour, wild honey, a sweet dish
  \item an arrow, a sword, armour, a flag, a drum
  \item a sheath, a treasure, a rupee, a stake, a fee
  \item news, a drama, a poem, a film, a game
\end{itemize}

\begin{tcolorbox}[breakable,colback=teal!4,colframe=teal!35!black,title={Before you move on}]
A tube? (\textbf{\sk{नालिका}}, \emph{nālikā}.) Buttermilk? (\textbf{\sk{तक्रम्}}, \emph{takram}.) Fruit juice? (\textbf{\sk{फलरसः}}, \emph{phalarasaḥ}.) A clove? (\textbf{\sk{लवंगम्}}, \emph{lavaṅgam}.) Cardamom? (\textbf{\sk{एला}}, \emph{elā}.) Puffed rice? (\textbf{\sk{लाजाः}}, \emph{lājāḥ}.) Grain? (\textbf{\sk{धान्यम्}}, \emph{dhānyam}.) Barley flour? (\textbf{\sk{सक्तुः}}, \emph{saktuḥ}.) Wild honey? (\textbf{\sk{क्षौद्रम्}}, \emph{kṣaudram}.) A sweet dish? (\textbf{\sk{मिष्टान्नम्}}, \emph{miṣṭānnam}.) An arrow? (\textbf{\sk{बाणः}}, \emph{bāṇaḥ}.) A sword? (\textbf{\sk{खड्गः}}, \emph{khaḍgaḥ}.) Armour? (\textbf{\sk{कवचम्}}, \emph{kavacam}.) A flag? (\textbf{\sk{ध्वजः}}, \emph{dhvajaḥ}.) A drum? (\textbf{\sk{पटहः}}, \emph{paṭahaḥ}.) A sheath? (\textbf{\sk{कोशः}}, \emph{kośaḥ}.) A treasure? (\textbf{\sk{निधिः}}, \emph{nidhiḥ}.) A rupee? (\textbf{\sk{रूप्यकम्}}, \emph{rūpyakam}.) A stake? (\textbf{\sk{पणः}}, \emph{paṇaḥ}.) A fee? (\textbf{\sk{शुल्कम्}}, \emph{śulkam}.) News? (\textbf{\sk{वार्ता}}, \emph{vārtā}.) A drama? (\textbf{\sk{नाटकम्}}, \emph{nāṭakam}.) A poem? (\textbf{\sk{काव्यम्}}, \emph{kāvyam}.) A film? (\textbf{\sk{चित्रपटः}}, \emph{citrapaṭaḥ}.) A game? (\textbf{\sk{खेला}}, \emph{khelā}.)
\end{tcolorbox}
